\section{引言}
\label{sec:introduction}

课程、编辑部或小型组织每天要判断大量文档是否值得进一步处理。人工审阅准确，但速度受限于
阅读与记录证据；纯自动分类快，却把「为什么这样判断」藏进模型内部。当决定会影响作者、
合作方或后续资源分配时，审阅者需要的不只是一个标签，而是一条可复核的证据链。

本文研究的问题是：\emph{在人机协同的文档分流中，能否在缩短审阅时间的同时，让人的决定和
支持它的证据保持可检查？} 我们以一个可复现的原型回答这个问题，而不是提出新的分类模型。

主要贡献有三点：
\begin{itemize}
  \item 一条三阶段的审阅工作流：文档进入、模型提出标签与证据、人工接受、修改或拒绝；
  \item 一个可复现的原型，包含固定环境、合成语料、审计记录与失败案例；
  \item 一次小规模内部评估，报告时间、推翻率与证据完整度，并明确其边界。
\end{itemize}

第~\ref{sec:related} 节讨论相关工作，第~\ref{sec:method} 节描述方法，
第~\ref{sec:implementation} 节说明实现与复现方式，第~\ref{sec:evaluation} 节报告评估，
第~\ref{sec:discussion} 节讨论伦理、责任与局限，第~\ref{sec:conclusion} 节总结。
